\centerline{\bf \Huge Вступительная проверочная работа}

\begin{flushright}
        \scriptsize{\textbf{Ответ Евклида египетскому царю Птолемею I, просившему указать ему более легкий путь изучения геометрии: — Нет царского пути в геометрии.}}
        
\end{flushright}

\problem{\textbf{1}} \textit{Аксиома; Теорема; Лемма; Следствие; Определение; Признак; Свойство}?

\begin{flushright} \textbf{1 балл}\\ 
\end{flushright}

\problem{\textbf{2}} \textit{Точка; Линия; Прямая линия; Луч; Отрезок; Параллельные и пересекающиеся прямые}?

\begin{flushright} \textbf{1 балл}\\ 
\end{flushright}

\problem{\textbf{3}} \textit{Угол; Вершина угла; Градус; Радиан; Развёрнутый, тупой, прямой и острый угол}?

\begin{flushright} \textbf{1 балл}\\ 
\end{flushright}

\problem{\textbf{4}} \textit{Смежные и вертикальные углы; Секущая; Свойства и признаки параллельных прямых и секущей}?

\begin{flushright} \textbf{1 балл}\\ 
\end{flushright}

\problem{\textbf{5}} \textit{Виды треугольников; Перпендикулярные прямые; Серединный перпендикуляр; Биссектриса угла; Биссектриса, медиана и высота треугольника}

\begin{flushright} \textbf{1 балл}\\ 
\end{flushright}

\problem{\textbf{6}} \textit{Свойства равных треугольников; Равнобедренный, равносторонний и прямоугольный треугольник}

\begin{flushright} \textbf{1 балл}\\ 
\end{flushright}

\problem{\textbf{7}} \textit{Параллелограмм, прямоугольник, ромб, квадрат, трапеция}

\begin{flushright} \textbf{3 балла}\\ 
\end{flushright}


\problem{\textbf{8}} \textit{Признаки и свойства параллелограмма, прямоугольника, ромба, квадрата, трапеции}

\begin{flushright} \textbf{3 балла}\\ 
\end{flushright}


\newpage
\hspace{3mm}



\problem{\textbf{1}} \textbf{Доказать:} 

\textit{Любой один признак параллельных прямых и секущей}

\begin{flushright} \textbf{3 балла}\\ 
\end{flushright}

\problem{\textbf{2}} \textbf{Доказать:} 

\textit{Признаки равенства треугольников (3б за каждый)}

\begin{flushright} \textbf{3 балла}\\ 
\end{flushright}



\problem{\textbf{3}} \textbf{Доказать:}

\textit{Если треугольники равны, то равны медианы, биссектрисы и высоты проведённые из соответствующих вершин (3б за каждый)}

\begin{flushright} \textbf{3 балла}\\ 
\end{flushright}



\problem{\textbf{4}} \textbf{Доказать:}

\textit{Признаки равенства прямоугольных треугольников (3б за каждый)}

\begin{flushright} \textbf{3 балла}\\ 
\end{flushright}

\problem{\textbf{5}} \textbf{Доказать:}

\textit{Признаки равнобедренного треугольника (3б за каждый)}

\begin{flushright} \textbf{3 балла}\\ 
\end{flushright}

\problem{\textbf{6}} \textbf{Доказать:}

\textit{Признак и свойство прямоугольного треугольника с углом в 30°(3б за каждый)} 

\begin{flushright} \textbf{3 балла}\\ 
\end{flushright}

\problem{\textbf{7}} \textbf{Доказать:}

\textit{Биссектриса и серединный перпендикуляр как ГМТ (3б за каждый)}

\begin{flushright} \textbf{3 балла}\\ 
\end{flushright}

\problem{\textbf{8}} \textbf{Вывод формулы площади треугольника} - \textbf{12 баллов}



